\clearpage
\section{Selecting Operations from Feedback}
After checking the budget, the rule calibrates if observed and believed support heights differ by more than 4 mm. Otherwise, failures under perturbation can trigger recovery demonstrations. If feedback SR is at least .85 and failures are diffuse with low perturbation pressure, the rule stops. If more than half the failures concentrate on one subgoal and fewer than 20,000 additional labels have been spent, it requests targeted regeneration. This last gate counts labels across the run, not labels on the failing subgoal: one 40,000-label update can therefore rule out targeted regeneration later. The rule also stops after two rounds without improvement; its default is a DAgger round. These decisions use deployment observations, \emph{not} the privileged expert contrast in Table~\ref{tab:probe}. A separate multi-label selector predicts geometry, coverage and recovery components from failure profiles; its training set excludes the three evaluation banks. Neither selector receives the injected cause at decision time.

We compare these choices with a fixed calibrate-once schedule, random and round-robin schedules, an upper-confidence-bound bandit (UCB) \cite{auer2002ucb}, a contextual linear bandit (LinUCB) \cite{linucb}, and a reference that knows the injected cause and plays a prescribed sequence. UCB explores operations based on observed improvement; LinUCB also conditions their estimated value on the failure profile. These are adaptations to the same operation menu, not published scores from the cited bandit papers.
\begin{table}[t]
\centering\caption{Selectors on four conditions. Entries are mean final SR over three banks and three training seeds (288 runs). First counts matches to the cause-informed reference's first operation, out of 36.}\label{tab:selectors}
\begin{tabular}{@{}lrrrrrr@{}}\toprule
Selector & Geom. & Cov. & Rec. & Healthy & Mean & First\\\midrule
Calibrate once & .772 & .600 & .833 & .911 & .779 & 9\\
Random & .300 & .939 & .867 & .972 & .769 & 0\\
Round robin & .894 & .872 & .822 & .939 & .882 & 9\\
UCB & .961 & .983 & .817 & .978 & .935 & 9\\
Rule & .900 & .928 & .900 & .928 & .914 & 21\\
LinUCB & .939 & .917 & .356 & .961 & .793 & 9\\
Multi-label & .828 & .928 & .828 & .883 & .867 & 28\\
Cause-informed & .939 & .972 & .939 & .844 & .924 & 36\\\bottomrule
\end{tabular}
\end{table}
The rule improved mean SR by 13.5 pp over calibrating once, with a paired 95\% percentile interval of $[8.6,18.6]$ pp (5,000 resamples of matched cell differences). The largest gain was on coverage: the rule could acquire missing layouts immediately. UCB had the highest observed mean. Its difference from the rule was $+2.1$ pp $[-1.1,5.1]$: it gained 5--6 points on geometry, coverage and healthy conditions, but lost 8.3 points on recovery. We do not treat the interval crossing zero as evidence of equivalence. The cause-informed schedule's corresponding difference was $+1.0$ pp $[-2.2,4.0]$; it is a reference for action agreement, not a maximum achievable SR.

\textbf{Success versus cost.} The multi-label selector matched more reference first actions than the rule (28 versus 21), yet its mean final SR was 4.7 points lower (paired interval $[-8.5,-1.0]$). On healthy tasks it used an average of 22,222 new labels, versus 44,444 for the rule and 80,000 for calibrate-once. The reference spent none: it stopped after observing feedback. To compare speed, we counted feedback episodes up to the first attainment of 90\% of the condition-specific reference mean. Multi-label averaged 29.3 episodes \emph{among the 27 of 36 runs that attained the threshold}; the rule averaged 35.6 among 33 of 36. A smaller conditional mean is not evidence that a selector reaches the threshold more reliably.

\textbf{More than one defect.} When height error and coverage deficit coexist, both a fixed schedule and a rule can start correctly by calibrating. Their second decisions can differ: the rule can now regenerate data on missing layouts. On a separate mixed-condition screen (four mixtures, two banks, one seed), UCB attained mean SR .863, the rule .731, calibrate-once .637 and the reference .900. Offline, a multi-label diagnosis recovered at least one true component for all 18 mixed profiles and all components for 72\%; single-label mixed-class recall was 6\%. These are different experiments. Detecting one component motivates repairing it and measuring the residual failure again; the offline result does not attribute UCB's return to a classifier it never used.
